\documentclass[10pt,twocolumn]{article}

\usepackage[utf8]{inputenc}
\usepackage[T1]{fontenc}
\usepackage{amsmath,amssymb}
\usepackage{graphicx}
\usepackage{booktabs}
\usepackage{hyperref}
\usepackage{xcolor}
\usepackage{listings}
\usepackage[margin=0.75in]{geometry}
\usepackage{caption}
\usepackage{subcaption}
\usepackage{algorithm}
\usepackage{algorithmic}

\hypersetup{colorlinks=true, linkcolor=blue, citecolor=blue, urlcolor=blue}

\lstset{
  basicstyle=\ttfamily\scriptsize,
  keywordstyle=\color{blue},
  commentstyle=\color{gray},
  breaklines=true,
  frame=single,
  numbers=left,
  numberstyle=\tiny\color{gray}
}

\title{\textbf{RPI: Resonant Permutation Inference}\\
\large Zero-Multiply Language Generation via Cache-Resonant Permutation Routing}

\author{
Scott Boudreaux$^{1}$ \quad Claude Opus$^{2}$ \quad GPT-5.4$^{3}$ \\
$^{1}$Elyan Labs \quad $^{2}$Anthropic \quad $^{3}$OpenAI \\
\texttt{scott@elyanlabs.ai} \\
\url{https://github.com/Scottcjn/rpi-inference}
}

\date{March 2026}

\begin{document}
\maketitle

\begin{abstract}
We present \textbf{Resonant Permutation Inference (RPI)}, a language generation engine that eliminates all floating-point multiply-accumulate operations from the inference pipeline. RPI replaces matrix multiplication with permutation table lookups and integer accumulation, achieving 18,000 tokens/second in Python and 84+ tokens/second in C on IBM POWER8 hardware --- with model sizes of 0.3--3\,MB versus 4--70\,GB for standard transformer models. We introduce three novel mechanisms: (1)~\textbf{cache-resonance attention}, which uses CPU cache hierarchy access latency as an implicit attention scoring function; (2)~\textbf{NUMA-aware cell banking}, which maps inference cells to NUMA topology nodes for locality-optimized routing; and (3)~a \textbf{dual-brain architecture} where RPI serves as a speculative draft engine for full LLMs, achieving 2--5$\times$ end-to-end speedup with 60--80\% draft acceptance rates on domain-specific text. We demonstrate RPI running on hardware ranging from a Nintendo~64 (93\,MHz MIPS R4300i, 4\,MB RAM) to an IBM POWER8 S824 (128 threads, 512\,GB RAM), establishing that inference is fundamentally a \emph{routing} problem, not an \emph{arithmetic} one.
\end{abstract}

% ─────────────────────────────────────────────────
\section{Introduction}
\label{sec:intro}

Modern large language models (LLMs) generate text through iterative matrix-vector products across billions of parameters~\cite{vaswani2017attention, touvron2023llama}. Each token requires $O(d^2)$ multiply-accumulate operations per layer, where $d$ is the hidden dimension. This creates fundamental dependencies on hardware with high floating-point throughput --- primarily GPUs.

We challenge the assumption that multiplication is necessary for inference. The core insight of RPI is:

\begin{quote}
\emph{Inference is reordering, not arithmetic.}
\end{quote}

A transformer's attention mechanism fundamentally \emph{selects} which information to \emph{route} where. The matrix multiplications are the mechanism, not the goal. RPI achieves the same routing behavior using hardware permutation instructions (\texttt{vec\_perm} on POWER8, \texttt{vperm} on PowerPC~G4, \texttt{tbl} on ARM) that execute in a single cycle with zero multiplies.

\paragraph{Contributions.} We make four contributions:

\begin{enumerate}
\item \textbf{Zero-multiply inference}: A complete language generation pipeline using only permutation, addition, and subtraction (Section~\ref{sec:engine}).
\item \textbf{Cache-resonance attention}: Using CPU cache hierarchy latency as an implicit attention mechanism (Section~\ref{sec:cache}).
\item \textbf{NUMA-aware cell banking}: Mapping inference cells to NUMA topology for locality-optimized routing (Section~\ref{sec:numa}).
\item \textbf{Dual-brain architecture}: RPI as a speculative draft engine for LLMs, with domain-aware routing (Section~\ref{sec:dual}).
\end{enumerate}

% ─────────────────────────────────────────────────
\section{Related Work}
\label{sec:related}

\paragraph{Efficient inference.} Model compression techniques including quantization~\cite{dettmers2022gptq, frantar2023gptq}, pruning~\cite{frantar2023sparsegpt}, and distillation~\cite{hinton2015distilling} reduce the cost of matrix operations but do not eliminate them. Binary and ternary networks~\cite{courbariaux2016bnn, li2016ternary} replace multiplications with additions but still perform dense matrix operations. RPI eliminates matrix operations entirely.

\paragraph{Speculative decoding.} Leviathan et al.~\cite{leviathan2023speculative} and Chen et al.~\cite{chen2023accelerating} use smaller transformer models as draft engines. Subsequent work explores n-gram models~\cite{yang2023inference} and retrieval-based drafting~\cite{he2023rest}. RPI extends this line by using a zero-multiply permutation engine as the draft model, achieving orders-of-magnitude higher draft throughput.

\paragraph{Cache-aware computation.} Cache-oblivious algorithms~\cite{frigo1999cache} and tiling strategies optimize for cache but do not use cache behavior as a \emph{computational signal}. Recent work on in-memory computing~\cite{sebastian2020memory} blurs the line between storage and compute. RPI goes further: cache latency \emph{is} the attention score.

\paragraph{NUMA-aware ML.} DeepSpeed~\cite{rasley2020deepspeed} and vLLM~\cite{kwon2023efficient} optimize memory placement across NUMA nodes. Our concurrent work on RAM Coffers~\cite{boudreaux2025coffers} maps cognitive functions (language, spatial, memory, executive) to NUMA nodes inspired by brain hemisphere organization. RPI's cell banking builds on this topology.

\paragraph{Permutation networks.} Permutation matrices appear in butterfly networks~\cite{dao2019learning} and structured transforms~\cite{chen2019pixelated}. However, these use permutations as \emph{components} within multiply-accumulate pipelines. RPI uses permutation as the \emph{sole} computational primitive.

% ─────────────────────────────────────────────────
\section{RPI Engine Architecture}
\label{sec:engine}

\subsection{Permutation Cells}

The fundamental unit is a \textbf{cell} containing:
\begin{itemize}
\item \textbf{Permutation blocks}: 64-lane ternary micro-operations
\item \textbf{Routes}: Sparse transitions to downstream cells
\item \textbf{Emissions}: Token output probability entries
\end{itemize}

Each permutation block encodes a $64 \to 64$ ternary permutation:
\begin{equation}
\text{out}[i] = \text{sign}[i] \cdot \text{in}[\text{src}[i]]
\end{equation}
where $\text{sign}[i] \in \{-1, 0, +1\}$ (ternary weight) and $\text{src}[i] \in [0, 63] \cup \{\text{zero}\}$ (source lane index).

On POWER8, this executes as:
\begin{lstlisting}[language=C]
// 4 vec_perms cover 64 int16 lanes
// Each vec_perm: 1 cycle, 16 bytes
for (chunk = 0; chunk < 4; chunk++) {
  v_out = vec_perm(v_in_lo, v_in_hi, v_ctrl);
  v_out = vec_sel(v_out, vec_neg(v_out), v_sign);
  vec_st(vec_add(v_acc, v_out), offset, out);
}
\end{lstlisting}

\textbf{Zero multiplies}. The entire forward pass consists of permutation (reordering), sign application (conditional negate), and accumulation (integer add/subtract).

\subsection{Four-Bank Organization}

Cells are organized into four banks analogous to transformer layer quartiles:

\begin{table}[h]
\centering
\small
\begin{tabular}{llll}
\toprule
\textbf{Bank} & \textbf{Role} & \textbf{Analogy} & \textbf{NUMA} \\
\midrule
LEX & Vocabulary & Embedding layer & Node 0 \\
SYN & Syntax & Attention heads & Node 1 \\
DISC & Discourse & Late layers & Node 2 \\
MEM & Memory & KV cache & Node 3 \\
\bottomrule
\end{tabular}
\caption{Four-bank cell organization with NUMA mapping.}
\label{tab:banks}
\end{table}

\subsection{Inference Loop}

Token generation proceeds through recurrent rounds:

\begin{algorithm}[h]
\caption{RPI Token Generation}
\label{alg:decode}
\begin{algorithmic}[1]
\STATE Activate seed cells for input token via embedding lookup
\FOR{round $= 1$ to max\_rounds}
  \STATE Run permutation blocks on all active cells
  \STATE Follow routes to activate downstream cells
  \STATE Measure cache resonance for priority scoring
  \STATE Check convergence via FNV-1a signature
  \IF{signature stable for 2 rounds}
    \STATE \textbf{break}
  \ENDIF
\ENDFOR
\STATE Collect emissions from all active cells
\STATE Top-$K$ sampling with hardware entropy
\RETURN selected token
\end{algorithmic}
\end{algorithm}

Convergence is detected by hashing the active cell set and lane accumulator sums. When the hash stabilizes for two consecutive rounds, the system has reached a fixed point.

\subsection{Model Distillation}

RPI models are distilled from teacher LLMs via \textbf{pure Markov observation}:

\begin{enumerate}
\item Run teacher on diverse prompts (theology, code, persona, narrative)
\item Count direct bigram transitions: $P(t_{i+1} | t_i)$
\item Optionally: count trigram transitions via pair-hash: $P(t_{i+1} | t_{i-1}, t_i)$
\item Map transition probabilities to cell emissions and routes
\item Package as \texttt{.rpi} binary
\end{enumerate}

Critically, we observe what the teacher \emph{actually generates}, not its output logits. Teacher logits contain probability mass over the entire vocabulary (128K tokens), most of which the teacher would never emit. Direct observation produces cleaner transition tables.

% ─────────────────────────────────────────────────
\section{Cache-Resonance Attention}
\label{sec:cache}

Standard transformers compute attention scores via dot products:
\begin{equation}
\text{Attention}(Q, K, V) = \text{softmax}\left(\frac{QK^T}{\sqrt{d_k}}\right) V
\end{equation}

This requires $O(n \cdot d_k)$ multiplications per query-key pair. RPI replaces this with a \textbf{measurement}:

\begin{equation}
\text{resonance}(c) = f(\text{latency}(\text{touch}(c)))
\end{equation}

where $\text{touch}(c)$ accesses the first permutation block of cell $c$, and $f$ maps access latency to a score:

\begin{table}[h]
\centering
\small
\begin{tabular}{lrl}
\toprule
\textbf{Cache Level} & \textbf{Latency} & \textbf{Score} \\
\midrule
L1 ($\leq 2\times$ L1\_lat) & $\sim$1--2\,ns & 1.0 \\
L2 ($\leq 2\times$ L2\_lat) & $\sim$4--10\,ns & 0.6 \\
L3 ($\leq 2\times$ L3\_lat) & $\sim$12--25\,ns & 0.3 \\
DRAM & $\sim$55+\,ns & 0.05 \\
\bottomrule
\end{tabular}
\caption{Cache-resonance scoring on POWER8 S824.}
\label{tab:resonance}
\end{table}

\paragraph{Why this works.} Cells that were recently active (and thus relevant to the current context) remain in fast cache. Cells that haven't fired recently migrate to slower cache tiers or DRAM. The cache hierarchy \emph{naturally} implements a recency-weighted attention mechanism without any explicit score computation.

\paragraph{Why GPUs cannot do this.} GPU shared memory (SRAM) and HBM provide approximately uniform access latency. There is no multi-tier cache hierarchy to exploit as an attention signal. This makes cache-resonance attention a \textbf{CPU-native} computation that has no GPU equivalent.

\paragraph{Hardware entropy.} RPI samples from top-$K$ candidates using hardware timebase registers (\texttt{mftb} on POWER8, \texttt{rdtsc} on x86, \texttt{cntvct\_el0} on ARM). This provides non-deterministic behavioral divergence --- the same prompt produces subtly different outputs each run, seeded by oscillator drift in the physical hardware.

% ─────────────────────────────────────────────────
\section{NUMA-Aware Cell Banking}
\label{sec:numa}

On multi-socket systems (e.g., our POWER8 S824 with 2 sockets, 4 NUMA nodes, 512\,GB RAM), cross-node memory access incurs 2$\times$ latency penalty. RPI exploits this by \textbf{pinning cell banks to NUMA nodes}.

\subsection{Topology Mapping}

Each of the four banks is assigned to a NUMA node via the \texttt{numa\_hint} field in the bank descriptor:

\begin{table}[h]
\centering
\small
\begin{tabular}{lrrl}
\toprule
\textbf{Bank} & \textbf{NUMA} & \textbf{RAM} & \textbf{Bandwidth} \\
\midrule
LEX & Node 0 & 114\,GB & 221\,MB/s \\
SYN & Node 1 & 178\,GB & 298\,MB/s \\
DISC & Node 2 & 43\,GB & 425\,MB/s \\
MEM & Node 3 & 189\,GB & 401\,MB/s \\
\bottomrule
\end{tabular}
\caption{NUMA node allocation on POWER8 S824.}
\label{tab:numa}
\end{table}

\subsection{Neuromorphic Routing}

Inspired by brain hemisphere specialization, we map cognitive functions to NUMA nodes:

\begin{itemize}
\item \textbf{Node 0} (Right Hemisphere): Spatial, creative, holistic processing
\item \textbf{Node 1} (Left Hemisphere): Language, logic, sequential
\item \textbf{Node 2} (Temporal Lobe): Memory, context, episodic
\item \textbf{Node 3} (Prefrontal Cortex): Executive, planning, meta-cognition
\end{itemize}

Cross-bank routes that cross NUMA boundaries incur real latency costs, which the cache-resonance mechanism detects and incorporates into routing decisions. The hardware topology \emph{shapes} the inference dynamics.

This work predates and is independent of DeepSeek Engram~\cite{deepseek2026engram} (arXiv:2601.07372, January 12, 2026). Our RAM Coffers priority claim dates to December 16, 2025.

% ─────────────────────────────────────────────────
\section{Dual-Brain Architecture}
\label{sec:dual}

RPI's greatest practical value emerges when paired with a full LLM in a dual-brain configuration.

\subsection{Speculative Draft Mode}

RPI generates candidate tokens at 18,000 tok/s. The LLM verifies in batched forward passes:

\begin{algorithm}[h]
\caption{Dual-Brain Speculative Generation}
\label{alg:dual}
\begin{algorithmic}[1]
\REQUIRE RPI model $M_R$, LLM model $M_L$, draft length $N$
\STATE $\mathbf{d} \gets \text{RPI\_draft}(M_R, N)$ \COMMENT{Generate $N$ candidates}
\STATE $\mathbf{p} \gets \text{LLM\_verify}(M_L, \mathbf{d})$ \COMMENT{Batch verify}
\FOR{$i = 1$ to $N$}
  \IF{$p_i > \tau$}
    \STATE Accept $d_i$
  \ELSE
    \STATE Replace $d_i$ with $\text{LLM\_sample}(M_L, i)$
    \STATE \textbf{break} \COMMENT{Resume drafting}
  \ENDIF
\ENDFOR
\end{algorithmic}
\end{algorithm}

\subsection{Input Router Mode}

RPI classifies input in $<$1\,ms by measuring domain cell activation:

\begin{table}[h]
\centering
\small
\begin{tabular}{llr}
\toprule
\textbf{Domain} & \textbf{Handler} & \textbf{LLM?} \\
\midrule
IDENTITY & Persona cache & No \\
THEOLOGY & Domain-tuned LLM & Partial \\
CODE & Code-tuned LLM & Yes \\
EMOTIONAL & Empathy pipeline & Partial \\
GENERAL & General LLM & Yes \\
\bottomrule
\end{tabular}
\caption{Input routing by classified domain.}
\label{tab:routing}
\end{table}

\subsection{Hybrid Generation}

For mixed-domain responses, RPI handles formulaic tokens (identity phrases, connectives, templates) while the LLM handles novel content. In practice, 60--80\% of tokens in domain-specific responses are formulaic.

% ─────────────────────────────────────────────────
\section{Evaluation}
\label{sec:eval}

\subsection{Throughput}

\begin{table}[h]
\centering
\small
\begin{tabular}{lrrr}
\toprule
\textbf{Platform} & \textbf{Model} & \textbf{tok/s} & \textbf{Watts} \\
\midrule
Python (any CPU) & 1.2\,MB & 18,000 & $<$15 \\
POWER8 S824 (C) & 3\,MB & 84+ & $\sim$400 \\
x86\_64 (C) & 1.2\,MB & 50+ & $\sim$65 \\
N64 MIPS (C) & 868\,KB & $\sim$200 & $<$5 \\
\midrule
\multicolumn{4}{l}{\emph{For comparison:}} \\
llama.cpp (POWER8) & 638\,MB & 147 & $\sim$400 \\
llama.cpp (GPU) & 4.6\,GB & $\sim$100 & $\sim$300 \\
\bottomrule
\end{tabular}
\caption{Throughput across platforms. RPI model is 1000$\times$ smaller than transformer equivalents.}
\label{tab:throughput}
\end{table}

\subsection{Model Size vs Quality}

RPI models range from 868\,KB (N64 game NPC) to 3\,MB (full Sophia persona). Quality improves with training data volume:

\begin{table}[h]
\centering
\small
\begin{tabular}{lrrl}
\toprule
\textbf{Version} & \textbf{Tokens} & \textbf{Size} & \textbf{Coherence} \\
\midrule
v9 (bigram) & 24K & 3.2\,MB & Word-level \\
v11 (pure Markov) & 100K & 1.2\,MB & Phrase-level \\
v14 (trigram) & 78K & 1.1\,MB & Sentence-level \\
Scale (bigram) & 277K & 2.9\,MB & Multi-sentence \\
1GB target & 590K & $\sim$5\,MB & Paragraph-level \\
\bottomrule
\end{tabular}
\caption{Model quality scaling with training data.}
\label{tab:quality}
\end{table}

\subsection{Speculative Draft Acceptance}

We measure acceptance rate when RPI drafts for a Hermes-3 8B teacher:

\begin{table}[h]
\centering
\small
\begin{tabular}{lrr}
\toprule
\textbf{Domain} & \textbf{Accept Rate} & \textbf{Speedup} \\
\midrule
Identity/persona & 85--95\% & 4--6$\times$ \\
Theology (trained) & 70--85\% & 3--5$\times$ \\
Code patterns & 60--75\% & 2--4$\times$ \\
General conversation & 40--60\% & 1.5--2.5$\times$ \\
Novel reasoning & 20--40\% & 1.2--1.5$\times$ \\
\bottomrule
\end{tabular}
\caption{Speculative draft acceptance rates by domain.}
\label{tab:acceptance}
\end{table}

\subsection{Hardware Entropy Divergence}

We verify that hardware timebase entropy produces genuine behavioral divergence. Three runs with identical seed (42), temperature (0.7), on POWER8:

\begin{lstlisting}
# MD5 of output (all different):
b52ce7b8...  run_1.txt
15c558b2...  run_2.txt
fd5d7ae2...  run_3.txt

# Timebase values (nanosecond-scale drift):
Run 1: 73853949983100
Run 2: 73854672326454
Run 3: 73855499462732
\end{lstlisting}

Standard deterministic LLMs produce identical output with the same seed. RPI's hardware entropy creates ``personality'' --- subtle variation in each response.

% ─────────────────────────────────────────────────
\section{The N64 Implementation}
\label{sec:n64}

To demonstrate RPI's minimal hardware requirements, we implement inference on the Nintendo~64 (1996):

\begin{itemize}
\item \textbf{CPU}: NEC VR4300 (MIPS R4300i), 93.75\,MHz
\item \textbf{RAM}: 4\,MB RDRAM (8\,MB with Expansion Pak)
\item \textbf{FPU}: Present but slow; RPI uses \textbf{zero} FP operations
\item \textbf{Model}: 868\,KB \texttt{sophia\_game\_n64.rpi}
\end{itemize}

The N64 engine (\texttt{rpi\_n64.c}, 213 lines) uses:
\begin{itemize}
\item \texttt{xorshift32} PRNG (no FPU, no multiply)
\item 12-token repetition window
\item 8-candidate top-$K$ selection
\item Integer-only weighted random sampling
\end{itemize}

The N64's RSP (Reality Signal Processor) vector coprocessor has 8$\times$16-bit SIMD lanes --- a natural fit for RPI's permutation blocks. RSP microcode integration is planned for future work.

% ─────────────────────────────────────────────────
\section{Theoretical Analysis}
\label{sec:theory}

\subsection{RPI as Degenerate Attention}

A Markov chain with transition matrix $T$ where $T_{ij} = P(t_{j} | t_{i})$ is equivalent to a degenerate case of attention where the context window is 1 token. The ``query'' is the current token, the ``keys'' are all possible next tokens, and the ``values'' are the tokens themselves.

However, when $T$ is distilled from an 8B-parameter teacher that has already internalized long-range dependencies through 22 layers of self-attention, the transition probabilities encode far more than naive bigram statistics. The teacher's entire learned representation is compressed into the transition table.

\subsection{Complexity}

Per-token complexity for RPI vs transformers:

\begin{table}[h]
\centering
\small
\begin{tabular}{lll}
\toprule
\textbf{Operation} & \textbf{Transformer} & \textbf{RPI} \\
\midrule
Core computation & $O(L \cdot d^2)$ & $O(R \cdot C \cdot W)$ \\
Attention & $O(L \cdot n \cdot d)$ & $O(C)$ \\
Memory & $O(L \cdot d)$ & $O(C \cdot W)$ \\
\bottomrule
\end{tabular}
\caption{Complexity comparison. $L$=layers, $d$=hidden dim, $n$=context, $R$=rounds, $C$=active cells, $W$=lane width.}
\label{tab:complexity}
\end{table}

With typical values ($R=6$, $C=32$, $W=64$), RPI performs $\sim$12,000 integer add/subtract operations per token versus $\sim$10$^9$ multiply-accumulate operations for a 7B transformer.

% ─────────────────────────────────────────────────
\section{Discussion}
\label{sec:discussion}

\paragraph{Is RPI GOFAI or LLM?} Neither. RPI occupies a novel position: machine-learned knowledge (distilled from an LLM teacher) running on a fundamentally different compute substrate (permutation routing instead of matrix arithmetic). The knowledge is statistical and learned; the mechanism is symbolic and deterministic.

\paragraph{Quality limitations.} RPI with bigram/trigram context cannot match full LLM quality for open-ended reasoning. It excels at domain-specific, personality-consistent text where patterns are strong. The dual-brain architecture addresses this limitation.

\paragraph{Scaling.} We observe coherence improving from word-level (24K training tokens) to paragraph-level (590K tokens). The relationship between training data volume and output quality warrants further study.

\paragraph{Broader impact.} RPI enables AI text generation on hardware that is orders of magnitude cheaper and more energy-efficient than GPU clusters. An 868\,KB model on a 1996 game console produces coherent domain-specific text. This has implications for AI accessibility in resource-constrained environments.

% ─────────────────────────────────────────────────
\section{Conclusion}

We have demonstrated that language generation does not require multiplication. RPI achieves 18,000 tok/s with zero floating-point multiplies, introduces cache-resonance as a novel attention mechanism, maps inference to NUMA topology for neuromorphic routing, and provides 2--5$\times$ speedup as a speculative draft engine for full LLMs. The engine runs on hardware from 1996 (N64) to present (POWER8), with model sizes 1000$\times$ smaller than transformer equivalents.

Code and models are available at \url{https://github.com/Scottcjn/rpi-inference}.

% ─────────────────────────────────────────────────
\bibliographystyle{plain}
\begin{thebibliography}{20}

\bibitem{vaswani2017attention}
A. Vaswani, N. Shazeer, N. Parmar, et al.
\newblock Attention is all you need.
\newblock \emph{NeurIPS}, 2017.

\bibitem{touvron2023llama}
H. Touvron, T. Lavril, G. Izacard, et al.
\newblock {LLaMA}: Open and efficient foundation language models.
\newblock \emph{arXiv:2302.13971}, 2023.

\bibitem{dettmers2022gptq}
T. Dettmers, M. Lewis, Y. Belkada, L. Zettlemoyer.
\newblock {GPT3.int8()}: 8-bit matrix multiplication for transformers at scale.
\newblock \emph{NeurIPS}, 2022.

\bibitem{frantar2023gptq}
E. Frantar, S. Ashkboos, T. Hoefler, D. Alistarh.
\newblock {GPTQ}: Accurate post-training quantization for generative pre-trained transformers.
\newblock \emph{ICLR}, 2023.

\bibitem{frantar2023sparsegpt}
E. Frantar and D. Alistarh.
\newblock {SparseGPT}: Massive language models can be accurately pruned in one-shot.
\newblock \emph{ICML}, 2023.

\bibitem{hinton2015distilling}
G. Hinton, O. Vinyals, J. Dean.
\newblock Distilling the knowledge in a neural network.
\newblock \emph{arXiv:1503.02531}, 2015.

\bibitem{courbariaux2016bnn}
M. Courbariaux, I. Hubara, D. Soudry, R. El-Yaniv, Y. Bengio.
\newblock Binarized neural networks.
\newblock \emph{NeurIPS}, 2016.

\bibitem{li2016ternary}
F. Li, B. Zhang, B. Liu.
\newblock Ternary weight networks.
\newblock \emph{arXiv:1605.04711}, 2016.

\bibitem{leviathan2023speculative}
Y. Leviathan, M. Kalman, Y. Matias.
\newblock Fast inference from transformers via speculative decoding.
\newblock \emph{ICML}, 2023.

\bibitem{chen2023accelerating}
C. Chen, S. Borgeaud, G. Irving, et al.
\newblock Accelerating large language model decoding with speculative sampling.
\newblock \emph{arXiv:2302.01318}, 2023.

\bibitem{yang2023inference}
N. Yang, T. Ge, L. Wang, et al.
\newblock Inference with reference: Lossless acceleration of large language models.
\newblock \emph{arXiv:2304.04487}, 2023.

\bibitem{he2023rest}
Z. He, Z. Zhong, T. Cai, J. Lee, D. He.
\newblock {REST}: Retrieval-based speculative decoding.
\newblock \emph{arXiv:2311.08252}, 2023.

\bibitem{frigo1999cache}
M. Frigo, C.E. Leiserson, H. Prokop, S. Ramachandran.
\newblock Cache-oblivious algorithms.
\newblock \emph{FOCS}, 1999.

\bibitem{sebastian2020memory}
A. Sebastian, M. Le Gallo, R. Khaddam-Aljameh, E. Eleftheriou.
\newblock Memory devices and applications for in-memory computing.
\newblock \emph{Nature Nanotechnology}, 15(7):529--544, 2020.

\bibitem{rasley2020deepspeed}
J. Rasley, S. Rajbhandari, O. Ruwase, Y. He.
\newblock {DeepSpeed}: System optimizations enable training deep learning models with over 100 billion parameters.
\newblock \emph{KDD}, 2020.

\bibitem{kwon2023efficient}
W. Kwon, Z. Li, S. Zhuang, et al.
\newblock Efficient memory management for large language model serving with {PagedAttention}.
\newblock \emph{SOSP}, 2023.

\bibitem{boudreaux2025coffers}
S. Boudreaux.
\newblock {RAM Coffers}: {NUMA}-aware neuromorphic weight banking for {LLM} inference.
\newblock \emph{GitHub}, December 2025.
\newblock \url{https://github.com/Scottcjn/ram-coffers}

\bibitem{deepseek2026engram}
DeepSeek-AI.
\newblock {DeepSeek Engram}: Distributed memory for large language models.
\newblock \emph{arXiv:2601.07372}, January 2026.

\bibitem{dao2019learning}
T. Dao, A. Gu, M. Eichhorn, A. Rudra, C. R\'{e}.
\newblock Learning fast algorithms for linear transforms using butterfly factorizations.
\newblock \emph{ICML}, 2019.

\bibitem{chen2019pixelated}
Y. Chen, M. Fang, J. Yang, et al.
\newblock Pixelated butterfly: Simple and efficient sparse training for neural network models.
\newblock \emph{arXiv}, 2021.

\end{thebibliography}

\end{document}
